\documentclass[12pt]{article}

\usepackage[utf8]{inputenc}
\usepackage[T1]{fontenc}
\usepackage{amsmath,amssymb,amsthm}
\usepackage{mathtools}
\usepackage{enumerate}
\usepackage[margin=2.5cm]{geometry}
\usepackage{hyperref}

% Theorem environments
\newtheorem{theorem}{Theorem}[section]
\newtheorem{proposition}[theorem]{Proposition}
\newtheorem{lemma}[theorem]{Lemma}
\newtheorem{corollary}[theorem]{Corollary}
\newtheorem{conjecture}[theorem]{Conjecture}
\theoremstyle{definition}
\newtheorem{definition}[theorem]{Definition}
\newtheorem{example}[theorem]{Example}
\theoremstyle{remark}
\newtheorem{remark}[theorem]{Remark}

% Shortcuts
\newcommand{\N}{\mathbb{N}}
\newcommand{\R}{\mathbb{R}}
\newcommand{\E}{\mathbb{E}}
\newcommand{\One}{\textsc{One}}
\newcommand{\All}{\textsc{All}}

\title{Optimal strategies in the all-heads coin game}
\author{Peter Pfaffelhuber}
\date{\today}

\begin{document}
\maketitle

\begin{abstract}
  We study a sequential coin-flipping game in which a player starts
  with~$n$ coins, each landing heads independently with
  probability~$p$.  In each round the player flips all remaining coins
  and must set aside at least one coin showing heads; if no coin shows
  heads, the player loses.  The player wins if and when all coins have
  been set aside.  This is a stochastic shortest-path Markov decision
  process whose Bellman equation involves a nonlinear suffix-maximum
  operator.  We analyse two natural strategies --- \One{} (set aside
  exactly one head) and \All{} (set aside every head) --- and
  determine the optimal winning probability~$w_{n,p}$ as a function
  of~$n$ and~$p$.

  For $p=\tfrac12$ every strategy that accepts the immediate win
  achieves $w_{n,1/2}=\tfrac12$.  For $p>\tfrac12$ the
  strategy~\One{} is optimal and $n\mapsto w_{n,p}$ is strictly
  increasing.  For $p<\tfrac12$ near~$\tfrac12$ we carry out a
  first-order perturbation expansion in $\delta:=\tfrac12-p$ and prove
  that the leading coefficient~$c_n$ of the deficit
  $\tfrac12-w_{n,p}$ satisfies a linear recursion for $n\ge 7$, with
  a limit $L=\lim c_n\approx 1.7035$.  As a consequence the optimal
  value sequence has, to first order in~$\delta$, a strict local
  minimum at $n=5$ and no local maximum --- local maxima are a
  non-perturbative phenomenon.
\end{abstract}

%% ====================================================================
\section{Introduction}\label{sec:intro}
%% ====================================================================

\subsection{The game}

Fix a probability $p\in(0,1)$ and a positive integer~$n$.  A player
is given $n$~coins, each of which independently lands heads with
probability~$p$ and tails with probability $q:=1-p$.  The game
proceeds in rounds.  In each round:
\begin{enumerate}[(i)]
\item All remaining coins are flipped simultaneously.
\item If no coin shows heads, the game is lost.
\item Otherwise the player must set aside at least one coin that
  shows heads.  The set-aside coins leave the game permanently.
\item If every coin has been set aside, the game is won.
  Otherwise, a new round begins with the remaining coins.
\end{enumerate}
The player's objective is to maximise the probability of winning.

Two natural strategies present themselves:
\begin{itemize}
\item \textbf{Strategy~\One:} set aside exactly one head when
  $1\le k\le n-1$ heads appear, and set aside all~$n$ when $k=n$.
\item \textbf{Strategy~\All:} set aside every head that appears
  ($i=k$).
\end{itemize}
Strategy~\One{} is conservative: it preserves as many coins as
possible for future rounds, betting that the one set-aside coin has
already ``done its job.''  Strategy~\All{} is greedy: it cashes in
every success immediately.  Neither is universally optimal.

A variant of this game, in which the objective is to maximise the
expected number of heads collected rather than the probability of
collecting all of them, was introduced and studied by van~Doorn
\cite{vanDoorn2024}.  The all-heads-wins formulation studied here
leads to a richer structural picture.

\subsection{Markov decision processes and the Bellman equation}

The game is a finite-state \emph{Markov decision process} (MDP): the
state is the number~$m$ of coins remaining, the action is the number
$i\in\{1,\dots,k\}$ of heads to set aside (where $k$ is the random
number of heads), and the transition is
$m\mapsto m-i$.  There are two absorbing states: \emph{win}
($m=0$, all coins collected) and \emph{lose} (reached when $k=0$ in
some round).

Because every policy absorbs with probability at least
$p^m+(1-p)^m>0$ per round, the MDP is a \emph{stochastic
shortest-path} (SSP) problem in the sense of Bertsekas and Tsitsiklis
\cite{BertsekisTsitsiklis1991}: the Bellman operator is a contraction
and the optimal value function is its unique fixed point.

\begin{definition}[Optimal winning probability]\label{def:w}
  Let $w_{0,p}:=1$.  For $n\ge 1$ define
  \begin{equation}\label{eq:bellman}
    w_{n,p} \;:=\; p^n + \sum_{j=1}^{n-1}\binom{n}{j}p^{n-j}q^{j}\,
    \max_{j\le m\le n-1} w_{m,p},
  \end{equation}
  where $q=1-p$ and the empty sum (at $n=1$) gives
  $w_{1,p}=p$.
\end{definition}

The term $p^n$ accounts for the event that all $n$~coins land heads
(an immediate win).  Each summand with $j\ge 1$ accounts for the
event that exactly~$j$ coins land tails; the player then chooses to
keep some number $m\in\{j,\dots,n-1\}$ of coins (setting aside
$n-m$ heads), and the game continues from state~$m$ with optimal
value~$w_{m,p}$.

The Bellman equation~\eqref{eq:bellman} is the standard optimality
equation for this SSP--MDP.  For general background on Markov decision
processes and the Bellman equation we refer to the monographs of
Puterman~\cite{Puterman1994}, Bertsekas~\cite{Bertsekas2017}, and
Stokey, Lucas and Prescott~\cite{StokeyLucas1989}.

\subsection{Strategies \One{} and \All}

Under strategy~\One{} the player always keeps $m=n-1$ coins
(except when $k=n$), so the Bellman recursion collapses to
\begin{equation}\label{eq:a-rec}
  a_{n,p} \;=\; p^n + \bigl(1-p^n-(1-p)^n\bigr)\,a_{n-1,p},
  \qquad a_{0,p}:=1.
\end{equation}
Under strategy~\All{} the player always keeps $m=j$ coins (the
number of tails), giving
\begin{equation}\label{eq:b-rec}
  b_{n,p} \;=\; p^n + \sum_{j=1}^{n-1}\binom{n}{j}p^{n-j}q^{j}\,
  b_{j,p},
  \qquad b_{0,p}:=1.
\end{equation}
Note that the recursion for~$b$ is \emph{linear} (no $\max$
operator), which makes it analytically more tractable.

\subsection{Outline}

In Section~\ref{sec:fair} we prove that at $p=\tfrac12$ the optimal
winning probability is $w_{n,1/2}=\tfrac12$ for every~$n$, and that
in fact every reasonable strategy achieves this value.  In
Section~\ref{sec:above} we show that for $p>\tfrac12$ the
strategy~\One{} is always optimal and $n\mapsto w_{n,p}$ is strictly
increasing.  Section~\ref{sec:below} treats the case $p<\tfrac12$
near~$\tfrac12$: we introduce a first-order perturbation in
$\delta=\tfrac12-p$, derive a recursion for the leading
coefficient~$c_n$, prove that it simplifies to a linear recursion for
$n\ge 7$, and compute the limit $L=\lim c_n$.

%% ====================================================================
\section{The fair coin: \texorpdfstring{$p=\tfrac12$}{p = 1/2}}
\label{sec:fair}
%% ====================================================================

\begin{theorem}\label{thm:half}
  For every $n\ge 1$, $w_{n,1/2}=\tfrac12$.  More generally, for any
  policy~$\pi$ satisfying $i_\pi(n,n)=n$ for every~$n$ (i.e.\ the
  player accepts the immediate win when all coins land heads), the
  winning probability is $v_n^\pi=\tfrac12$.

  In particular, both strategies~\One{} and~\All{} achieve the
  optimum: $a_{n,1/2}=b_{n,1/2}=w_{n,1/2}=\tfrac12$.
\end{theorem}

\begin{proof}
  We use the binomial identity
  \begin{equation}\label{eq:bin-half}
    \sum_{j=1}^{n-1}\binom{n}{j}\Bigl(\frac{1}{2}\Bigr)^{n}
    \;=\; 1-2^{1-n}.
  \end{equation}

  \emph{Step~1.}  We show $w_{n,1/2}=\tfrac12$ by strong induction.
  For $n=1$: $w_{1,1/2}=\tfrac12$.  For $n\ge 2$, suppose
  $w_{m,1/2}=\tfrac12$ for all $1\le m\le n-1$.  Then every
  suffix-maximum equals~$\tfrac12$, so
  \begin{align*}
    w_{n,1/2}
    &= 2^{-n} + \frac{1}{2}\sum_{j=1}^{n-1}\binom{n}{j}2^{-n} \\
    &= 2^{-n} + \tfrac{1}{2}\bigl(1-2^{1-n}\bigr) \\
    &= 2^{-n} + \tfrac{1}{2} - 2^{-n}
    \;=\; \tfrac{1}{2}.
  \end{align*}

  \emph{Step~2.}  For a general policy~$\pi$ with
  $i_\pi(n,n)=n$, conditioning on the number of heads
  $K\sim\mathrm{Bin}(n,\tfrac12)$ gives
  \[
    v_n^\pi = 2^{-n}\cdot 1 + \sum_{k=1}^{n-1}\binom{n}{k}2^{-n}\,
    v_{n-i_\pi(n,k)}^\pi.
  \]
  Since $1\le i_\pi(n,k)\le k\le n-1$ for $k<n$, the subscript
  $n-i_\pi(n,k)$ lies in $\{1,\dots,n-1\}$, and by strong induction
  each $v_{n-i_\pi(n,k)}^\pi=\tfrac12$.  The same calculation as in
  Step~1 gives $v_n^\pi=\tfrac12$.
\end{proof}

\begin{remark}\label{rem:refuse}
  A policy that refuses the immediate win at $k=n$ (i.e.\
  $i_\pi(n,n)<n$) achieves strictly less than~$\tfrac12$: the
  contribution of the $k=n$ event drops from $2^{-n}$ to
  $2^{-n}\cdot v_{n-i_\pi(n,n)}^\pi < 2^{-n}$, while all other terms
  remain at most $\tfrac12\cdot\binom{n}{k}2^{-n}$, giving
  $v_n^\pi<\tfrac12$.
\end{remark}

%% ====================================================================
\section{Above the fair coin: \texorpdfstring{$p>\tfrac12$}{p > 1/2}}
\label{sec:above}
%% ====================================================================

\begin{theorem}\label{thm:above}
  For $p>\tfrac12$ and all $n\ge 2$:
  \begin{enumerate}[\rm(i)]
  \item $w_{n,p}>w_{n-1,p}$, i.e.\ $n\mapsto w_{n,p}$ is strictly
    increasing.
  \item $w_{n-1,p}<\dfrac{p^n}{p^n+q^n}$.
  \item Strategy~\One{} is optimal: $w_{n,p}=a_{n,p}$.
  \end{enumerate}
\end{theorem}

\begin{proof}
  We prove (i) and (ii) simultaneously by strong induction on~$n$.

  \emph{Base case $n=2$.}
  One computes $w_{1,p}=p$ and
  $w_{2,p}=p^2+(1-p^2-q^2)p = 3p^2-2p^3$.  Then
  $w_{2,p}-w_{1,p}=p(1-p)(2p-1)>0$ since $p>\tfrac12$.
  Also $w_{1,p}=p<p^2/(p^2+q^2)$ iff $p^2+q^2>p$, which holds for
  $p\in(\tfrac12,1)$.

  \emph{Inductive step.}  Assume (i) and (ii) hold for all indices up
  to~$n-1$.  Since $n\mapsto w_{n,p}$ is increasing on
  $\{1,\dots,n-1\}$ by hypothesis, the suffix-maximum
  $\max_{j\le m\le n-1}w_{m,p}=w_{n-1,p}$ for every
  $j\in\{1,\dots,n-1\}$.  Hence the Bellman
  equation~\eqref{eq:bellman} reduces to
  \[
    w_{n,p} = p^n + (1-p^n-q^n)\,w_{n-1,p},
  \]
  which is exactly the recursion~\eqref{eq:a-rec} for
  strategy~\One.  This proves~(iii).

  For~(i), $w_{n,p}>w_{n-1,p}$ iff
  $p^n(1-w_{n-1,p})>q^n w_{n-1,p}$, i.e.\
  $w_{n-1,p}<p^n/(p^n+q^n)$, which is hypothesis~(ii) at level
  $n-1$.

  For~(ii) at level~$n$, we need $w_{n,p}<p^{n+1}/(p^{n+1}+q^{n+1})$.
  Using the recursion and the inductive bound on~$w_{n-1,p}$, a
  direct computation shows this reduces to the inequality
  $(q/p)^n<(q/p)^{n-1}$, which holds since $q/p<1$.
\end{proof}

%% ====================================================================
\section{Below the fair coin: perturbation in
  \texorpdfstring{$\delta=\tfrac12-p$}{delta = 1/2 - p}}
\label{sec:below}
%% ====================================================================

For $p<\tfrac12$ the structure of $n\mapsto w_{n,p}$ is far richer:
numerical experiments show that the sequence develops local minima and
maxima, and the optimal strategy is neither~\One{} nor~\All{} in
general; see \cite{vanDoorn2024}, Section~10.

To get analytical traction, we study the regime
$p=\tfrac12-\delta$ with $\delta>0$ small.

\subsection{The deficit and its recursion}

\begin{definition}\label{def:deficit}
  For $p\in(0,1)$ and $n\ge 1$, define the
  \emph{deficit}
  $\Delta_{n,p}:=\tfrac12-w_{n,p}$.
\end{definition}

By Theorem~\ref{thm:half}, $\Delta_{n,1/2}=0$ for all~$n$.

\begin{proposition}\label{prop:delta-rec}
  The deficit satisfies
  \begin{equation}\label{eq:delta-rec}
    \Delta_{n,p} \;=\;
    \frac{q^n-p^n}{2}
    + \sum_{j=1}^{n-1}\binom{n}{j}p^{n-j}q^{j}\,
    \min_{j\le m\le n-1}\Delta_{m,p},
  \end{equation}
  with $\Delta_{0,p}=-\tfrac12$.
  In particular:
  \begin{enumerate}[\rm(i)]
  \item For $p<\tfrac12$: $\Delta_{n,p}>0$ for every $n\ge 1$
    (i.e.\ $w_{n,p}<\tfrac12$).
  \item For $p>\tfrac12$: $\Delta_{n,p}<0$ for every $n\ge 1$
    (i.e.\ $w_{n,p}>\tfrac12$).
  \end{enumerate}
\end{proposition}

\begin{proof}
  Substitute $w_{m,p}=\tfrac12-\Delta_{m,p}$ into the
  Bellman equation~\eqref{eq:bellman}.  The $\max$ on
  $w_m$ becomes a $\min$ on~$\Delta_m$.  The constant terms
  $\tfrac12\cdot\bigl[\text{binomial weights}\bigr]$ cancel by the
  identity $p^n+\sum_{j=1}^{n-1}\binom{n}{j}p^{n-j}q^j\cdot\tfrac12
  = \tfrac12 - \tfrac12(q^n-p^n)$, leaving~\eqref{eq:delta-rec}.
  The sign claims follow by induction using $q^n-p^n>0$ (resp.\ $<0$)
  and the positivity (resp.\ negativity) of the $\min$ terms.
\end{proof}

\subsection{First-order expansion}

Set $p=\tfrac12-\delta$ with $\delta>0$ small.  Then
$q=\tfrac12+\delta$, and the constant term in~\eqref{eq:delta-rec} is
\[
  \frac{q^n-p^n}{2}
  = \frac{(\tfrac12+\delta)^n-(\tfrac12-\delta)^n}{2}
  = \frac{n\delta}{2^{n-1}} + O(\delta^3).
\]
The binomial weights become $\binom{n}{j}2^{-n}+O(\delta)$, and the
$\min$ terms are $O(\delta)$ by Proposition~\ref{prop:delta-rec}.
Thus $\Delta_{n,p}=O(\delta)$.

\begin{proposition}[First-order coefficient]\label{prop:cn}
  Write $\Delta_{n,p}=c_n\delta+O(\delta^2)$ as $\delta\to 0^+$.
  Then $c_1=1$ and for $n\ge 2$:
  \begin{equation}\label{eq:cn-rec}
    c_n \;=\; \frac{n}{2^{n-1}}
    + \frac{1}{2^n}\sum_{j=1}^{n-1}\binom{n}{j}\,
    \min_{j\le m\le n-1}c_m.
  \end{equation}
\end{proposition}

\begin{proof}
  Match first-order coefficients in~\eqref{eq:delta-rec}.  The
  constant term contributes $n/2^{n-1}$.  Each binomial weight is
  $\binom{n}{j}2^{-n}+O(\delta)$; the $\min$ of the
  $\Delta_m=c_m\delta+O(\delta^2)$ at first order is
  $(\min c_m)\delta+O(\delta^2)$.  Products of $O(\delta)$ terms are
  $O(\delta^2)$.
\end{proof}

\begin{example}\label{ex:cn-small}
  The first few values of~$c_n$:
  \[
    c_1=1,\quad c_2=\tfrac{3}{2},\quad c_3=\tfrac{27}{16},\quad
    c_4=\tfrac{111}{64},\quad
    c_5=\tfrac{3555}{2048}\approx 1.7358.
  \]
  The sequence increases on $\{1,\dots,5\}$ and decreases for
  $n\ge 6$.
\end{example}

\subsection{Collapse of the minimum}

The key structural insight is that for $n\ge 7$, the $\min$ in
\eqref{eq:cn-rec} can be evaluated explicitly.

\begin{lemma}[Asymptotic collapse]\label{lem:collapse}
  For all $n\ge 7$ and $j\in\{1,\dots,n-1\}$:
  \begin{equation}\label{eq:collapse}
    \min_{j\le m\le n-1}c_m
    \;=\;
    \begin{cases}
      c_j     & \text{if } j\in\{1,2,3\},\\
      c_{n-1} & \text{if } j\in\{4,\dots,n-1\}.
    \end{cases}
  \end{equation}
\end{lemma}

The proof of Lemma~\ref{lem:collapse} requires two auxiliary claims.

\begin{lemma}\label{lem:lower}
  $c_m\ge\tfrac{27}{16}$ for every $m\ge 4$.
\end{lemma}

\begin{proof}
  We use the algebraic identity
  \begin{equation}\label{eq:alg-id}
    A_n - \tfrac{27}{16}B_n
    \;=\; -\frac{3(n^2-15n+36)}{32\cdot 2^n},
  \end{equation}
  where $A_n$ and $B_n$ are defined in~\eqref{eq:AB} below.

  \emph{For $m\in\{4,\dots,12\}$:} the quadratic $n^2-15n+36$ has
  roots at $n=3$ and $n=12$, so it is non-positive on
  $\{3,\dots,12\}$.  Hence $A_m-\tfrac{27}{16}B_m\ge 0$ for these
  values, and the linear recursion~\eqref{eq:cn-linear} gives
  $c_m\ge\tfrac{27}{16}$ by induction.

  \emph{For $m\ge 13$:} define the deficit
  $\varepsilon_m:=c_m-\tfrac{27}{16}$.
  From the linear recursion,
  \[
    \varepsilon_m \;=\; (1-B_m)\varepsilon_{m-1}
    + \bigl(A_m-\tfrac{27}{16}B_m\bigr).
  \]
  The starting value $\varepsilon_{12}=c_{12}-\tfrac{27}{16}>\tfrac{1}{60}$
  (verified by exact rational computation).
  The cumulative negative contribution is bounded by
  $\sum_{k\ge 13}\tfrac{3(k^2-15k+36)}{32\cdot 2^k}
  \le\tfrac{3}{32}\cdot\tfrac{99}{2048}=\tfrac{297}{65536}<\tfrac{1}{200}$,
  using $\sum_{k\ge 13}k^2/2^k=\tfrac{99}{2048}$.  Meanwhile
  $\prod_{m\ge 13}(1-B_m)>\tfrac{7}{8}$ (by
  $\sum_{m\ge 13}B_m=\tfrac{243}{2048}<\tfrac{1}{8}$).
  So
  $\varepsilon_m>\tfrac{1}{60}\cdot\tfrac{7}{8}-\tfrac{1}{200}
  =\tfrac{23}{2400}>0$ for all $m\ge 13$.
\end{proof}

\begin{lemma}\label{lem:decreasing}
  The sequence $(c_n)_{n\ge 5}$ is strictly decreasing.
\end{lemma}

\begin{proof}
  For $n\ge 7$, the linear recursion gives
  $c_n-c_{n-1}=A_n-B_nc_{n-1}$.  By Lemma~\ref{lem:lower},
  $c_{n-1}\ge\tfrac{27}{16}$, so
  $c_n-c_{n-1}\le A_n-\tfrac{27}{16}B_n<0$ for $n\ge 13$
  (by~\eqref{eq:alg-id} with $n^2-15n+36>0$).  For
  $n\in\{5,\dots,12\}$ one checks $c_n<c_{n-1}$ by direct
  computation.
\end{proof}

\begin{proof}[Proof of Lemma~\ref{lem:collapse}]
  For $j\in\{1,2,3\}$: since $c_1<c_2<c_3<\tfrac{27}{16}\le c_m$ for
  $m\ge 4$ (Lemma~\ref{lem:lower}), the minimum over
  $\{j,\dots,n-1\}$ is $c_j$.

  For $j\ge 4$: all values $c_j,\dots,c_{n-1}$ are at
  least~$\tfrac{27}{16}$, and the sequence is decreasing for
  indices~$\ge 5$ (Lemma~\ref{lem:decreasing}).  For $j=4$
  one checks $c_4>\min(c_5,\dots,c_{n-1})=c_{n-1}$ since
  $c_4=\tfrac{111}{64}\approx 1.734>c_6\approx 1.729$.
  Hence the minimum is $c_{n-1}$.
\end{proof}

\subsection{Linear recursion and the limit}

\begin{proposition}[Linear recursion]\label{prop:linear}
  For $n\ge 7$:
  \begin{equation}\label{eq:cn-linear}
    c_n \;=\; A_n + (1-B_n)\,c_{n-1},
  \end{equation}
  where
  \begin{equation}\label{eq:AB}
    A_n := \frac{n}{2^{n-1}}
    + \frac{n\,c_1+\binom{n}{2}c_2+\binom{n}{3}c_3}{2^n},
    \qquad
    B_n := \frac{2+n+\binom{n}{2}+\binom{n}{3}}{2^n}.
  \end{equation}
\end{proposition}

\begin{proof}
  Substitute \eqref{eq:collapse} into \eqref{eq:cn-rec} and
  split the sum at $j=3$.  For $j\in\{1,2,3\}$ the $\min$ is
  $c_j$; for $j\ge 4$ the $\min$ is $c_{n-1}$.  Collecting terms
  gives~\eqref{eq:cn-linear}--\eqref{eq:AB}.
\end{proof}

\begin{theorem}[Existence and value of the limit]\label{thm:limit}
  The limit $L:=\lim_{n\to\infty}c_n$ exists, and for any $n_0\ge 7$:
  \begin{equation}\label{eq:L-formula}
    L \;=\; c_{n_0-1}\prod_{m\ge n_0}(1-B_m)
    \;+\; \sum_{k\ge n_0}A_k\prod_{m>k}(1-B_m).
  \end{equation}
  Both the infinite product and the series converge at geometric rate
  since $A_n,B_n=O(n^3/2^n)$.  Numerically,
  \begin{equation}\label{eq:L-value}
    L \;=\; 1.70347176087173673645\ldots
  \end{equation}
\end{theorem}

\begin{proof}
  From \eqref{eq:cn-linear}, $c_n=A_n+(1-B_n)c_{n-1}$.  Since
  $B_n>0$ and $A_n>0$ for all~$n$, and $(c_n)$ is eventually
  decreasing and bounded below (by Lemma~\ref{lem:lower}), the limit
  exists.  Iterating the recursion from $n=n_0$ and taking
  $n\to\infty$ yields~\eqref{eq:L-formula}; the rate follows from
  $B_n=O(n^3/2^n)$.
\end{proof}

\subsection{Shape of the optimal-value sequence near
  \texorpdfstring{$p=\tfrac12$}{p=1/2}}

\begin{corollary}\label{cor:shape}
  For $\delta>0$ sufficiently small and $p=\tfrac12-\delta$:
  \begin{enumerate}[\rm(i)]
  \item $w_{n-1,p}-w_{n,p}=(c_n-c_{n-1})\delta+O(\delta^2)$.
  \item The sequence $n\mapsto w_{n,p}$ has, to first order
    in~$\delta$, a strict local minimum at $n=5$ (since $c_5>c_4$
    and $c_5>c_6$).
  \item There is no local maximum at first order: $c_n$ is eventually
    decreasing, so $w_{n,p}$ is eventually increasing.
  \end{enumerate}
\end{corollary}

\begin{remark}
  Numerical computation shows that the sequence $n\mapsto w_{n,p}$
  does develop strict local maxima for $p$ bounded away
  from~$\tfrac12$ (e.g.\ at $n=9$ for $p=0.42$, verified by exact
  rational arithmetic).  This is consistent with the corollary: local
  maxima are a \emph{non-perturbative} phenomenon invisible to the
  first-order expansion.  Van Doorn~\cite{vanDoorn2024}, Section~10,
  conjectures that the optimal strategy for $p<\tfrac12$ is
  ``infinitely complex.''
\end{remark}

%% ====================================================================
%\section*{Acknowledgements}
%% ====================================================================

\begin{thebibliography}{99}

\bibitem{Bertsekas2017}
  D.~P.~Bertsekas,
  \emph{Dynamic Programming and Optimal Control},
  Vol.~I, 4th ed.,
  Athena Scientific, Belmont, MA, 2017.

\bibitem{BertsekisTsitsiklis1991}
  D.~P.~Bertsekas and J.~N.~Tsitsiklis,
  An analysis of stochastic shortest path problems,
  \emph{Math.\ Oper.\ Res.}\ \textbf{16} (1991), 580--595.

\bibitem{DubinsSavage1965}
  L.~E.~Dubins and L.~J.~Savage,
  \emph{How to Gamble If You Must: Inequalities for Stochastic
    Processes},
  McGraw-Hill, New York, 1965;
  reprinted by Dover, 2014.

\bibitem{Puterman1994}
  M.~L.~Puterman,
  \emph{Markov Decision Processes: Discrete Stochastic Dynamic
    Programming},
  Wiley, New York, 1994.

\bibitem{StokeyLucas1989}
  N.~L.~Stokey, R.~E.~Lucas~Jr., and E.~C.~Prescott,
  \emph{Recursive Methods in Economic Dynamics},
  Harvard University Press, Cambridge, MA, 1989.

\bibitem{vanDoorn2024}
  E.~A.~van~Doorn,
  Optimal play of a coin-flipping game,
  \emph{arXiv preprint} arXiv:2406.14700, 2024.

\end{thebibliography}

\end{document}
